% The three parts of the Day 15 lab. The steps follow Labs/day15/README.md.
% The rubric and the requirements are the text of Docs/capstone.md and of
% Section 7 of the syllabus.

% ------------------------------------------------------------------ Part A
\begin{frame}{Part A: build and measure (80 min)}
  \small
  \begin{enumerate}
    \item \textbf{Complete the system (40 min).} Connect the parts. Test
          each requirement of your proposal once.
    \item \textbf{Measure (30 min).} Fill the benchmark table of
          \code{report.md}. Measure the size and the accuracy of the model
          before and after the optimization on the same test set.
    \item \textbf{Rehearse (10 min).} Run the demonstration once with the
          clock: the live system, the cut of the network link, the
          dashboard. Prepare plan B: a recorded input or a short video.
  \end{enumerate}
  \vskip 0.4em
  \begin{block}{Before Part A}
    The lecture time (100 minutes) is build time. Build the parts that the
    demonstration needs first. Follow the plan of your Day 14 report.
  \end{block}
\end{frame}

\begin{frame}{Part A: the benchmark table}
  \footnotesize
  \renewcommand{\arraystretch}{1.15}
  \setlength{\tabcolsep}{3pt}
  \begingroup\centering
  \begin{tabular}{@{}L{1.5cm}L{1.3cm}L{0.9cm}L{1.3cm}L{1.0cm}L{1.6cm}L{1.6cm}@{}}
    \toprule
    \textbf{Model} & \textbf{Board} & \textbf{Format} & \textbf{Median}
      & \textbf{p95} & \textbf{Memory} & \textbf{Energy} \\
    \midrule
    \ldots & XIAO & \code{int8} & \ldots\ ms & \ldots\ ms & arena
      \ldots\ bytes & \ldots\ mJ \\
    \ldots & Pi 5 & \ldots & \ldots\ ms & \ldots\ ms & RSS \ldots\ MB &
      \ldots\ mJ \\
    \bottomrule
  \end{tabular}\par\endgroup
  \vskip 0.5em
  \small
  \textbf{The method of Day 9}, in the report under the table:
  \begin{itemize}
    \item warm-up runs and the number of timed runs,
    \item the window: the model only, or end to end,
    \item the statistic: the median and a high percentile,
    \item the energy: a USB power meter, or the power values of Day 9
          (say ``estimate'').
  \end{itemize}
  \source{The form of the table in \texttt{Labs/day15/report.md}. The
    values come from your boards.}
\end{frame}

\begin{frame}{Part A: the five requirements, a last check}
  \footnotesize
  \renewcommand{\arraystretch}{1.15}
  \setlength{\tabcolsep}{4pt}
  \begingroup\centering
  \begin{tabular}{@{}lL{0.42\textwidth}L{0.44\textwidth}@{}}
    \toprule
    \textbf{No.} & \textbf{Requirement} & \textbf{Show it in the
      demonstration} \\
    \midrule
    1 & XIAOML Kit and Raspberry Pi & Each board with its own task \\
    2 & One optimized model, size and accuracy before and after & The
      table of Section 4 of the report \\
    3 & One local decision, no internet & Cut the link to the instructor
      broker; the decision continues \\
    4 & Telemetry with MQTT to a dashboard & The live dashboard \\
    5 & Benchmark table: latency, memory, energy & The table of the
      previous frame \\
    \bottomrule
  \end{tabular}\par\endgroup
  \vskip 0.6em
  \small
  A requirement that does not work: show what works, use plan B, and write
  the limit in the report.
\end{frame}

% ------------------------------------------------------------------ Part B
\begin{frame}{Part B: the demonstration (10 min)}
  \begin{columns}[T]
    \begin{column}{0.50\textwidth}
      \footnotesize
      \renewcommand{\arraystretch}{1.2}
      \setlength{\tabcolsep}{4pt}
      \begin{tabular}{@{}lL{3.9cm}@{}}
        \toprule
        \textbf{Time} & \textbf{Content} \\
        \midrule
        1 min & The problem and the requirements \\
        4 min & The system works: live, also with the network link cut \\
        3 min & The benchmark table, two decisions with their numbers \\
        2 min & Questions \\
        \bottomrule
      \end{tabular}
    \end{column}
    \begin{column}{0.46\textwidth}
      \small
      \begin{itemize}
        \item Start on time: the next team waits.
        \item Show, then explain. The live system comes before the slides.
        \item Answer with numbers.
        \item Watch the other teams. Write one question for each team.
      \end{itemize}
    \end{column}
  \end{columns}
  \vskip 0.6em
  \small
  Part B has 70 minutes: 7 teams with 10 minutes each. With more teams, the
  instructor uses the schedule of Day 14.
\end{frame}

\begin{frame}{Part B: the rubric}
  \footnotesize
  \renewcommand{\arraystretch}{1.15}
  \setlength{\tabcolsep}{4pt}
  \begingroup\centering
  \begin{tabular}{@{}L{0.30\textwidth}cL{0.52\textwidth}@{}}
    \toprule
    \textbf{Criterion} & \textbf{Weight} & \textbf{Excellent} \\
    \midrule
    The system works in the demonstration & 30\,\% & All five
      requirements work live, also with the link cut \\
    Integration of the two boards, messages, and monitoring & 25\,\% & Each
      board has a clear task; topics and sequence numbers; a live
      dashboard \\
    Measured results and analysis of trade-offs & 25\,\% & A complete
      table with the method; each main decision with two measured numbers
      and what you gave up \\
    Report and presentation & 20\,\% & Short, clear, the template; answers
      with numbers \\
    \bottomrule
  \end{tabular}\par\endgroup
  \vskip 0.5em
  \small
  \code{Docs/capstone.md}, Section 3, gives the levels ``adequate'' and
  ``insufficient''.
  \source{Weights: Section 7 of the syllabus. Levels of the third
    criterion: Decision Log rubric of the textbook instructor guide.}
\end{frame}

% ------------------------------------------------------------------ Part C
\begin{frame}{Part C: report and feedback (30 min)}
  \begin{columns}[T]
    \begin{column}{0.52\textwidth}
      \begin{block}{The report: at most four pages}
        \small
        \begin{enumerate}
          \item Problem and requirements
          \item The system
          \item The five requirements
          \item The optimized model
          \item The benchmark table
          \item Robustness
          \item Decision Log (two or three)
          \item Limits and next steps
        \end{enumerate}
      \end{block}
    \end{column}
    \begin{column}{0.44\textwidth}
      \small
      \begin{itemize}
        \item \textbf{Report (20 min):} complete \code{report.md}, and
              submit it as the instructor says.
        \item \textbf{Feedback (10 min):} fill \code{feedback.md} with no
              name: eight ratings, the level of each week, four open
              questions.
      \end{itemize}
    \end{column}
  \end{columns}
  \vskip 0.5em
  \small
  The example report of the machine monitor is in
  \code{solutions/report\_example.md}.
\end{frame}
